\RequirePackage{fix-cm}
\documentclass[twocolumn]{svjour3}
\smartqed
\usepackage{graphicx}
\usepackage{amsmath}
\usepackage{amssymb}
\usepackage{url}
\usepackage{microtype}
\emergencystretch=2em
\journalname{Journal of Real-Time Image Processing}

\newcommand{\blankfield}[1]{\textit{#1: }\underline{\hspace{3.0cm}}}

\begin{document}

\title{When TensorRT Is Not the Bottleneck: End-to-End PatchCore Optimization on Jetson Orin Nano}
\titlerunning{End-to-End PatchCore Optimization on Jetson Orin Nano}
\author{Minwoo Lee \and HyeongTae Ahn}
\authorrunning{Lee and Ahn}
\institute{Minwoo Lee \at Department of Computer Engineering, Kumoh National
Institute of Technology, Gumi, Republic of Korea
\and HyeongTae Ahn (corresponding author) \at Department of Computer
Engineering, Kumoh National Institute of Technology, Gumi, Republic of Korea \\
\blankfield{E-mail}}
\date{}
\maketitle

\begin{abstract}
Memory-bank anomaly detectors can retain strong unsupervised detection
accuracy, yet their deployment cost on edge hardware is easily understated
when evaluation is restricted to the neural feature extractor. We implement a
hash-traced PatchCore runtime in C++17, TensorRT, and CUDA on an NVIDIA Jetson
Orin Nano and measure the complete inference path, including preprocessing,
data transfers, feature extraction, embedding layout conversion, exact
nearest-neighbor retrieval, postprocessing, power, and energy. On the MVTec AD
bottle category, FP16 reduces feature-extraction time from 20.212 to 3.785 ms,
but end-to-end latency decreases by only 1.48\% because retrieval remains the
dominant stage. Under a fixed data path, asynchronous tiled CUDA search reduces
retrieval time by 35.70$\times$ and end-to-end latency by 28.43$\times$
relative to CPU search. A staged system ablation achieves a steady-state
interval of 482.350 ms, whereas a five-ratio coreset study reduces latency from
992.520 to 72.093 ms and bank storage from 250.799 to 12.539 MB while preserving
an image-level AUROC of 1.0. Across all 15 MVTec AD categories, the selected
runtime achieves macro image- and pixel-level AUROCs of 0.983314 and 0.979761,
respectively. In a 4,000-frame run lasting 33.94 min, it sustains 1.964
images/s, with a 517.845 ms p95 interval, a maximum GPU temperature of
63.906$^\circ$C, and no degradation beyond a predefined 5\% bound. These
results demonstrate that real-time claims for edge anomaly detection must be
grounded in complete-pipeline measurements rather than backbone timing alone.
\keywords{industrial anomaly detection \and PatchCore \and edge inference \and TensorRT \and CUDA \and real-time image processing}
\end{abstract}

\section{Introduction}\label{sec:introduction}
Industrial visual inspection commonly provides abundant defect-free samples
but few examples spanning the range of possible failures. Unsupervised anomaly
detection addresses this imbalance by modeling nominal appearance and treating
departures from that model as candidate defects. MVTec AD supports systematic
evaluation across 15 object and texture categories using image-level labels and
pixel-level masks \cite{bergmann2019mvtec}. PatchCore is representative of a
high-performing class of memory-bank methods: it embeds local image patches
with a pretrained network, retains a coreset of nominal embeddings, and assigns
anomaly scores from distances to that bank \cite{roth2022patchcore}.

The same memory bank that supports detection accuracy complicates edge
deployment. Accelerating the convolutional backbone does not reduce the cost
of comparing each test patch with thousands of high-dimensional nominal
vectors. On a resource-constrained platform, retrieval-kernel organization,
memory layout, transfer policy, buffer lifetime, and pipeline scheduling can
therefore determine system behavior after feature extraction has been
accelerated. Reporting TensorRT engine time alone may consequently obscure
end-to-end latency, energy consumption, and whether both throughput and
response-time requirements are met.

This work evaluates the complete PatchCore inference path on an NVIDIA Jetson
Orin Nano. It makes a systems contribution rather than proposing a new anomaly
score. Specifically, it provides:
\begin{enumerate}
\item a hash-traced C++/TensorRT/CUDA runtime that measures stage-wise latency,
accuracy, power, energy, memory, and artifact lineage across the full pipeline;
\item controlled ablations of numerical precision, coreset ratio, exact
retrieval backend, and the S0--S6 data path, including negative INT8 and
transpose results;
\item category-wise validation over all 15 MVTec AD categories using
independently trained 10\% coresets; and
\item a sustained thermal and timing evaluation that separates pipeline interval
from per-frame completion latency and identifies the evidence still needed for
a fixed-rate deadline claim.
\end{enumerate}

The results show that feature extraction can be reduced to a few milliseconds
while exact retrieval still requires hundreds of milliseconds. For the tested
bottle workload, the optimized system has p95 pipeline interval and completion
latency below reference bounds of 550 ms and 1.1 s. The repeated-frame trace
does not impose fixed-rate arrivals and therefore does not establish an
all-frames deadline guarantee.

\section{Related work}\label{sec:related}
\subsection{Industrial anomaly detection}
PatchCore constructs a representative coreset of nominal patch embeddings and
performs nearest-neighbor scoring in a pretrained feature space
\cite{roth2022patchcore}. Alternative families change the statistical model or
training objective. DRAEM jointly learns reconstruction and discrimination
from synthetically corrupted images \cite{zavrtanik2021draem}, whereas CFLOW-AD
uses conditional normalizing flows to estimate feature likelihoods
\cite{gudovskiy2022cflow}. N-Pad augments anomaly detection with positional and
local-neighborhood information \cite{jang2023npad}. These approaches advance
detection methodology; in contrast, the present study fixes the PatchCore
representation and asks how its inference graph behaves after deployment on a
resource-constrained device.

\subsection{Efficient and edge-oriented inference}
Model redesign can avoid or reduce the cost structure of a large memory bank.
EfficientAD, for example, adopts a lightweight student--teacher architecture
and reports millisecond-scale latency on a modern GPU
\cite{batzner2024efficientad}. PaSTe explicitly targets efficient visual
anomaly detection at the edge \cite{barusco2025paste}, while tiled ensembles
address the memory demands of high-resolution inspection
\cite{rolih2024tiled}. Such methods are complementary rather than direct
baselines for the present systems question. Once a memory-bank detector has
been selected, its achievable service rate depends on the execution of both
feature extraction and retrieval. A fair cross-model comparison would
therefore require a common device, input resolution, power mode, and
end-to-end measurement boundary.

\subsection{Similarity search and embedded GPU execution}
Large-scale GPU similarity search has been studied extensively in datacenter
settings \cite{johnson2021faissgpu}. Jetson-class systems, however, combine GPU
compute with constrained memory capacity and power, making data movement and
integration with the feature extractor part of the deployment problem. CUDA
provides streams, pinned memory, and asynchronous data-transfer mechanisms
\cite{nvidia2026cuda}, while TensorRT supports optimized lower-precision neural
inference \cite{nvidia2026tensorrt}. We use these mechanisms to isolate the
costs of exact PatchCore retrieval on an embedded platform. Vendor
documentation informs the implementation only; every accuracy and performance
claim reported below is derived from the measured experimental artifacts.

\section{Runtime design}\label{sec:runtime}
\subsection{PatchCore execution path}
For each $256\times256$ RGB input, WideResNet50-2 features from two intermediate
levels are spatially aligned and concatenated into
$E\in\mathbb{R}^{1536\times32\times32}$. Flattening the spatial axes produces
$Q\in\mathbb{R}^{1024\times1536}$. Given a nominal bank
$B\in\mathbb{R}^{M\times1536}$, exact patch scores are
\begin{equation}
d_i=\min_{1\leq j\leq M}\lVert Q_i-B_j\rVert_2^2.
\end{equation}
The largest reweighted patch distance determines the image score. The patch
scores are reshaped, interpolated, and Gaussian-smoothed to produce the anomaly
map. The 10\% bottle bank contains $M=21{,}401$ embeddings and occupies
125.396 MB in FP32.

Figure~\ref{fig:pipeline} summarizes the measured path. SHA-256 values link
each output to its executable, configuration, engine, memory bank, source
image, export manifest, and evaluation dataset. Accuracy and performance rows
are associated only when their executable, configuration, engine, and bank
hashes match.

\begin{figure*}[t]
\centering
\includegraphics[width=0.88\textwidth]{Fig1.pdf}
\caption{Measured inference path. The reported end-to-end time includes every block, while TensorRT time covers feature extraction only.}
\label{fig:pipeline}
\end{figure*}

\subsection{Exact CUDA retrieval}
The direct kernel assigns threads to query--bank comparisons, with limited
reduction efficiency and data reuse. The warp variant cooperatively reduces
distances within each warp. The tiled kernel stages fragments of bank vectors
in shared memory, and the asynchronous variant overlaps tile movement with
arithmetic where supported. All variants accumulate distances in FP32 and
resolve equal distances by selecting the lowest bank index, ensuring
deterministic tie handling.

TensorRT emits channel-major NCHW data, whereas distance evaluation benefits
from patch-major queries. A transpose-assisted backend consumes NCHW output on
device for the GPU-resident S5 path. We retain its S0 result as a negative
control because S0 already passes through host patch-major conversion.

\subsection{System stages and two deadlines}
S0 uses transient device buffers and pageable synchronous copies. S1 retains
device buffers; S2 adds pinned host staging; S3 uses asynchronous copies on the
default stream; S4 introduces a dedicated non-blocking stream; S5 keeps the
TensorRT output device-resident through retrieval; and S6 overlaps next-frame
preprocessing with current-frame runtime work.

S6 requires two distinct timing quantities. The pipeline interval between
consecutive completed outputs determines the sustainable arrival rate, whereas
completion time measures a frame's residence time in the two-frame pipeline.
We therefore evaluate the system against both an input-cycle bound and a
response-time bound.

\section{Experimental method}\label{sec:method}
\subsection{Platform and protocol}
Table~\ref{tab:platform} specifies the fixed hardware and software environment.
The Release executable targets compute capability 8.7. Unless noted otherwise,
accuracy is evaluated on the complete test split, whereas performance is
measured with one fixed normal image, 50 warm-up iterations, 200 measured
iterations, and three independent runs. The collected statistics include
stage-wise means, per-run p95 values, between-run variation, power, and energy;
telemetry is sampled at a nominal interval of 100 ms. The test split is not
used to select a threshold, coreset ratio, or numerical precision.

\begin{table}[t]
\caption{Experimental platform.}\label{tab:platform}
\centering\small
\begin{tabular}{ll}\hline
Device & Jetson Orin Nano, 8 GB\\
Power mode & MAXN\_SUPER\_ID2\\
JetPack / L4T & 6.2.1 / 36.4.7\\
CUDA / TensorRT & 12.6 / 10.3.0.30\\
OpenCV / compiler & 4.8.0 / GCC 11.4\\
Input / embedding & $256^2$ / $1536\times32\times32$\\
\hline
\end{tabular}
\end{table}

Before benchmarking, the implementation was checked against exported PyTorch
references at the preprocessed-input, embedding, patch, nearest-neighbor,
score, and anomaly-map levels. The FP32 TensorRT engine used for this reference
check was built with TF32 disabled. CPU and CUDA retrieval outputs agree within
the recorded numerical tolerances.

\subsection{Controlled ablations}
The precision experiment fixes the 10\% FP32 bank and S5 retrieval path while
varying the feature engine among FP32, FP16, and entropy-calibrated INT8.
Calibration uses 100 normal training images. The bank-storage experiment
deterministically serializes FP16 values from the FP32 export. For the coreset
experiment, independent bottle models are trained with seed 42 at ratios of 1,
2.5, 5, 10, and 20\%; the smaller banks are therefore not truncations of a
single export.

The retrieval experiment fixes FP16, the 10\% FP32 bank, S0, the executable,
and the power mode, varying only the six exact-search backends. The system
experiment fixes the transpose-assisted backend and varies S0--S6. The
multi-category evaluation uses one independently trained seed-42 10\% bank for
each category, the complete test split of that category, and three performance
runs.

\subsection{Sustained criterion}
The final FP16/S6 path is evaluated over 4,000 measured bottle frames. Before
analysis, operational degradation is defined as an increase of more than 5\%
in mean pipeline interval between the first and last complete five-minute
windows. The evaluated reference bounds are a 550 ms pipeline interval and a
1.1 s completion time. The repeated-frame run captures observed timing and
temperature behavior, but does not record fixed-rate arrivals, p99, maximum,
deadline misses, queue growth, or dropped frames. The available
\texttt{tegrastats} output does not provide a direct hardware-throttling flag.

\section{Results}\label{sec:results}
\subsection{Precision and coreset trade-offs}
As shown in Table~\ref{tab:precision}, FP16 reduces TensorRT time by 81.3\%,
but reduces end-to-end latency by only 1.48\%. Exact retrieval continues to
require 467.057 ms. Relative to FP16, INT8 changes total latency by only 4.94\%
while reducing image AUROC to 0.688889. We therefore reject this
entropy-calibrated post-training quantization configuration. This negative
result is specific to the evaluated configuration and does not establish the
behavior of quantization-aware training or selective mixed precision.

\begin{table}[t]
\caption{Controlled precision results on bottle.}\label{tab:precision}
\centering\scriptsize
\begin{tabular}{lrrrr}\hline
 & Image & Pixel & TRT & Total\\
Precision & AUROC & AUROC & (ms) & (ms)\\\hline
FP32 & 1.0000 & 0.9855 & 20.212 & 558.511\\
FP16 & 1.0000 & 0.9855 & 3.785 & 550.245\\
INT8 & 0.6889 & 0.7615 & 2.703 & 523.076\\
\hline
\end{tabular}
\end{table}

Serializing the bank in FP16 halves its stored payload, but runtime bank memory
remains 125.396 MB because the current loader expands the values to FP32. No
supported runtime speedup is observed from serialization alone.
Figure~\ref{fig:tradeoffs} also summarizes the coreset experiment. All five
ratios preserve an image AUROC of 1.0, while pixel AUROC ranges from 0.984961
to 0.985500. Reducing the ratio from 20\% to 1\% lowers mean latency from
992.520 to 72.093 ms, bank storage from 250.799 to 12.539 MB, and energy from
13.599 to 0.865 J/image. This Pareto relationship is established only for the
bottle category and seed 42.

\begin{figure*}[t]
\centering
\includegraphics[width=0.96\textwidth]{Fig2.pdf}
\caption{Precision and coreset trade-offs. INT8 loses substantial accuracy, whereas coreset size controls the dominant retrieval cost.}
\label{fig:tradeoffs}
\end{figure*}

\subsection{Retrieval backend ablation}
All six backends preserve image- and pixel-level AUROCs of 1.0 and 0.985452,
respectively. Table~\ref{tab:backend} and Fig.~\ref{fig:backend} show that direct
CUDA search is 3.28$\times$ faster end to end than CPU search. Warp reduction,
tiling, and asynchronous tiling progressively reduce total latency to 844.294,
613.777, and 510.285 ms. The best S0 backend provides a 35.70$\times$ retrieval
speedup and a 28.43$\times$ end-to-end speedup over CPU, while reducing energy
by 92.2\%.

\begin{table}[t]
\caption{Exact retrieval backend ablation (three runs).}\label{tab:backend}
\centering\scriptsize
\begin{tabular}{lrrr}\hline
Backend & NN ms & Total ms & J/image\\\hline
CPU & 14434.802 & 14505.975 & 94.681\\
Direct CUDA & 4314.865 & 4420.934 & 70.795\\
Warp & 735.978 & 844.294 & 9.541\\
Tiled & 508.859 & 613.777 & 9.943\\
Async tiled & 404.378 & 510.285 & 7.367\\
Async tiled + transpose & 422.996 & 527.955 & 7.772\\
\hline
\end{tabular}
\end{table}

Adding transpose assistance makes S0 3.46\% slower than asynchronous tiling
alone. This result is consistent with the evaluated data path: S0 already
performs patch-major conversion on the host, whereas transpose assistance is
intended for device-resident NCHW output.

\begin{figure*}[t]
\centering
\includegraphics[width=0.9\textwidth]{Fig3.pdf}
\caption{Controlled nearest-neighbor backend latency. A logarithmic axis exposes the large CPU and direct-CUDA gap.}
\label{fig:backend}
\end{figure*}

\subsection{System-stage ablation}
Every S0--S6 configuration retains image- and pixel-level AUROCs of 1.0 and
0.985455, respectively. Table~\ref{tab:system} indicates that persistent
allocation alone has little effect. Pinned staging at S2 produces the largest
isolated reduction in single-frame latency before S5. Neither S4 nor S5 lowers
the mean latency with the already optimized retrieval kernel; consequently,
these CUDA mechanisms should not be interpreted as independently sufficient
optimizations.

\begin{table}[t]
\caption{S0--S6 system ablation.}\label{tab:system}
\centering\scriptsize
\begin{tabular}{lrrr}\hline
Stage & Completion ms & Interval ms & J/image\\\hline
S0 & 548.333 & 548.333 & 7.121\\
S1 & 547.476 & 547.476 & 7.111\\
S2 & 502.427 & 502.427 & 6.577\\
S3 & 501.316 & 501.316 & 6.526\\
S4 & 508.594 & 508.594 & 7.195\\
S5 & 503.152 & 503.152 & 7.462\\
S6 & 964.625 & 482.350 & 6.426\\
\hline
\end{tabular}
\end{table}

Relative to S0, S6 reduces the steady-state interval by 12.0\% and energy by
9.8\%. Its completion time increases to 964.625 ms because two frames are in
flight. The 482.350 ms value is therefore a throughput interval rather than a
per-frame response latency (Fig.~\ref{fig:system}).

\begin{figure*}[t]
\centering
\includegraphics[width=0.88\textwidth]{Fig4.pdf}
\caption{System-stage completion time and pipeline interval. S6 improves service interval while increasing individual-frame completion time.}
\label{fig:system}
\end{figure*}

\subsection{Multi-category generality}
Across all 15 categories, the selected FP16/S5 runtime achieves macro image-
and pixel-level AUROCs of 0.983314 and 0.979761, respectively
(Fig.~\ref{fig:categories}). The corresponding image/pixel macro pairs are
0.978531/0.984692 for objects and 0.992882/0.969898 for textures. Toothbrush
has the lowest image AUROC (0.919444), while wood has the lowest pixel AUROC
(0.929994). These category-level results qualify any aggregate claim of strong
accuracy.

\begin{figure*}[t]
\centering
\includegraphics[width=0.98\textwidth]{Fig5.pdf}
\caption{Category-wise runtime accuracy and mean end-to-end latency for independently trained 10\% coresets.}
\label{fig:categories}
\end{figure*}

Mean category latency ranges from 164.050 ms for toothbrush (6,144 bank
entries) to 906.637 ms for hazelnut (40,038 entries). The macro mean latency is
574.334 ms, with mean power of 14.043 W and mean energy of 8.128 J/image.
TensorRT time remains within approximately 3.3--3.8 ms, whereas retrieval time
ranges from 119.257 to 780.257 ms. Retrieval is therefore the dominant cost
beyond the bottle category as well.

\subsection{Sustained thermal and deadline behavior}
The 4,000-frame run lasts 2,036.670 s. The mean pipeline interval is 509.168 ms
and its p95 is 517.845 ms, corresponding to 1.964 images/s. Mean completion
time is 1,018.309 ms, with a p95 of 1,031.660 ms. Both p95 values are below
the evaluated 550 ms interval and 1.1 s completion reference bounds. Because
the run does not model timestamped fixed-rate arrivals or record p99, maximum,
miss rate, backlog, and drops, it does not show that every frame meets a
response deadline.

The mean interval is 514.355 ms in the first complete five-minute window and
507.128 ms in the last, a change of $-1.405$\% that remains within the
predefined 5\% degradation criterion. The maximum GPU temperature is
63.906$^\circ$C, and mean/peak input power is 14.014/15.939 W.
Figure~\ref{fig:sustained} shows no progressive slowdown. A single 101\% GPU
utilization sample is treated as a telemetry artifact rather than evidence of
utilization above 100\%. Because no direct throttling flag is available, the
inference is limited to stability in the observed performance and temperature
measurements.

\begin{figure*}[t]
\centering
\includegraphics[width=0.9\textwidth]{Fig6.pdf}
\caption{Five-minute sustained-run windows. Mean and p95 interval remain stable as GPU temperature reaches steady state.}
\label{fig:sustained}
\end{figure*}

\section{Discussion}\label{sec:discussion}
\subsection{Interpreting the system bottleneck}
The experiments distinguish four quantities that answer different deployment
questions. TensorRT engine time measures feature extraction, and retrieval time
measures the memory-bank operator. Completion time characterizes a frame's
response, whereas the pipeline interval characterizes the sustainable service
rate when work is overlapped. The precision results show why backbone timing
cannot substitute for end-to-end timing: an 81.3\% reduction in TensorRT time
produces only a 1.48\% reduction in total latency. Conversely, the S6 pipeline
improves the steady-state interval while increasing completion time. An
optimization can therefore improve throughput without improving response
latency, and reporting only one of these metrics can reverse the apparent
conclusion.

This distinction also bounds the timing claim. The sustained experiment shows
p95 values below 550 ms interval and 1.1 s completion references for the tested
FP16/S6 bottle workload in the specified power mode. Fixed-rate deadline
compliance remains untested. The experiment also does not demonstrate
that every MVTec AD category meets those bounds: the S5 category means already
span 164.050--906.637 ms. Nor does the absence of observed slowdown establish
the absence of hardware throttling, because the available telemetry lacks a
direct throttle indicator.

\subsection{Implications for memory-bank deployment}
Two complementary levers govern exact retrieval. Kernel organization reduces
the cost of each comparison, as shown by the progression from direct CUDA to
warp, tiled, and asynchronous tiled search. Coreset reduction changes both the
number of comparisons and the amount of stored state. In the bottle experiment,
the 1\% bank dominates the tested larger banks with respect to image AUROC,
latency, storage, and energy. Its pixel AUROC nevertheless differs slightly
from those of the larger banks, and the comparison uses a single seed; the
result should therefore be understood as a measured bottle-specific frontier,
not a general prescription to use a 1\% coreset.

The multi-category results reinforce this interpretation. Feature-extraction
time remains within a narrow few-millisecond range, while retrieval varies by
more than a factor of six across category-specific banks. Accordingly, a fixed
backbone benchmark provides little information about the service rate of a
memory-bank detector unless bank cardinality and retrieval execution are also
reported. The evidence further suggests that deployment choices should be
evaluated jointly against accuracy, latency, memory, and energy rather than
optimized as independent objectives.

\subsection{Negative results and design boundaries}
Several plausible optimizations fail under the tested conditions. Entropy-
calibrated INT8 substantially reduces accuracy without eliminating the
retrieval bottleneck. FP16 bank serialization reduces the file payload but not
runtime memory because the loader expands the bank to FP32. Transpose
assistance regresses under S0, where host-side conversion has already produced
a patch-major layout, and a dedicated stream does not by itself reduce mean
latency. These findings show that the benefit of a mechanism depends on the
surrounding data path; isolated use of lower precision, asynchronous execution,
or a specialized layout is not sufficient evidence of an end-to-end gain.

\subsection{Limitations and future work}
The study is limited to one embedded platform, MVTec AD, one backbone, custom
exact brute-force retrieval, and a single training seed for each coreset condition.
The power and timing results are specific to the stated device configuration,
and the soft real-time envelope is specific to the bottle workload. The
sustained test establishes stability over 33.94 min rather than indefinite
operation, and the telemetry supports no direct conclusion about hardware
throttling. These constraints limit both statistical and platform-level
generalization.

Future evaluation requires matched cuBLAS and, where feasible on Jetson ARM64,
FAISS GPU exact-search baselines; validation-based multi-seed bank selection;
a native FP16 device bank with FP32 accumulation; and fixed-rate arrivals with
p99/max, misses, backlog, and drops. Additional industrial datasets, backbones,
and edge platforms should then be examined. Comparisons with model-efficient detectors should be conducted
on the same device and power mode, with matched input size and a common
end-to-end measurement boundary, so that architectural efficiency is not
confounded with experimental scope.

\section{Conclusion}\label{sec:conclusion}
Deploying PatchCore on the Jetson Orin Nano is governed by the complete
memory-bank pipeline rather than by TensorRT in isolation. FP16 reduces feature
extraction to a few milliseconds but leaves exact retrieval dominant. Under
the controlled S0 data path, CUDA kernel organization yields a 28.43$\times$
end-to-end speedup over CPU search; coreset size then governs much of the
remaining latency, storage, and energy. The selected runtime preserves strong
aggregate accuracy across all 15 MVTec AD categories, while the 33.94 min
repeated-frame bottle experiment shows stable p95 interval and completion
latency without exceeding the predefined degradation bound. Fixed-rate
deadline compliance remains to be tested. Complete-
pipeline measurement, explicit deadline definitions, and traceable artifacts
are therefore essential to support bounded real-time claims for memory-bank
anomaly detection at the edge.

\section*{Statements and Declarations}
\paragraph{Funding} \blankfield{Statement}

\paragraph{Competing interests} \blankfield{Statement}

\paragraph{Author contributions} \blankfield{Statement}

\paragraph{Data availability} MVTec AD is available from its original provider under its terms. Experiment-derived CSV files, manifests, configuration files, and analysis scripts supporting this article are available in the project repository at \url{https://github.com/minwoo1119/edge-anomaly-detection}. Large third-party dataset images and platform-specific TensorRT engines are not redistributed.

\paragraph{Code availability} Runtime, experiment orchestration, validation, and paper-result generation code is available at \url{https://github.com/minwoo1119/edge-anomaly-detection}. The frozen evidence manifest records SHA-256 values for the reported artifacts.

\paragraph{Use of generative AI} OpenAI Codex assisted with code inspection, experiment orchestration, result checking, and manuscript drafting. The authors are responsible for verifying the implementation, measurements, citations, interpretation, and final text.

\bibliographystyle{spmpsci}
\bibliography{references}

\end{document}
